\documentclass[10pt]{beamer}
\usetheme{Madrid}
\usepackage[T1]{fontenc}
\usepackage[utf8]{inputenc}
\usepackage[brazil]{babel}
\usepackage{lmodern}
\usepackage{hyperref}
\usepackage{listings}
\usepackage{xcolor}
\setbeamertemplate{navigation symbols}{}

\lstset{
  language=C++,
  basicstyle=\ttfamily\small,
  keywordstyle=\color{blue},
  commentstyle=\color{gray},
  stringstyle=\color{red},
  breaklines=true,
  showstringspaces=false,
}

\title{Tipos e Variáveis em C++}
\subtitle{Curso de Microcontroladores}
\author{}
\date{}

\begin{document}

\begin{frame}
  \titlepage
\end{frame}

\begin{frame}[fragile]{Tipos de dados básicos}
\begin{lstlisting}
int contador = 0;      // inteiro
float temperatura = 25.5f;  // número com decimal
char letra = 'A';      // um caractere
bool ligado = true;    // verdadeiro/falso
\end{lstlisting}

  \begin{itemize}
    \item `int`: inteiros, de -32768 a 32767 (em Arduino).
    \item `float`: números decimais, menos preciso que `double`.
    \item `char`: um único caractere entre apóstrofos.
    \item `bool`: apenas `true` ou `false`.
  \end{itemize}
\end{frame}

\begin{frame}[fragile]{Constantes em Arduino}
\begin{lstlisting}
const int LED_PIN = 13;
const float PI = 3.14159f;
const char LETRA_A = 'A';

void setup() {
  pinMode(LED_PIN, OUTPUT);
  // LED_PIN não pode mudar depois
}
\end{lstlisting}

  \begin{itemize}
    \item `const` impede que o valor mude depois.
    \item Maior clareza no código.
    \item Economiza memória (armazenado em Flash, não em RAM).
  \end{itemize}
\end{frame}

\begin{frame}[fragile]{Tipos otimizados}
\begin{lstlisting}
byte sensorValue = 0;    // 0-255
unsigned int leitura = analogRead(A0);
long tempo = millis();   // inteiro grande
\end{lstlisting}

  \begin{itemize}
    \item `byte`: 0 a 255, usa menos memória que `int`.
    \item `unsigned int`: positivos apenas.
    \item `long`: inteiro muito grande para valores de tempo.
    \item Em sistemas embarcados, otimizar tipo é essencial.
  \end{itemize}
\end{frame}

\begin{frame}[fragile]{Variáveis globais e locais}
\begin{lstlisting}
int contadorGlobal = 0;  // acessível em todo sketch

void setup() {
  int valorLocal = 10;   // só dentro de setup()
  contadorGlobal = 100;
}

void loop() {
  // valorLocal não existe aqui!
  // mas contadorGlobal sim
}
\end{lstlisting}

  \begin{itemize}
    \item Globais: declaradas fora de funções, acessíveis em qualquer lugar.
    \item Locais: dentro de funções, desaparecem quando função termina.
  \end{itemize}
\end{frame}

\begin{frame}[fragile]{Exemplo prático: sensor de temperatura}
\begin{lstlisting}
const int TEMP_PIN = A1;
float temperatura = 0;
bool alerta = false;

void loop() {
  int leitura = analogRead(TEMP_PIN);
  // Converter leitura (0-1023) para graus
  temperatura = (leitura / 1023.0) * 100;
  
  if (temperatura > 30.0) {
    alerta = true;
    digitalWrite(LED_PIN, HIGH);
  } else {
    alerta = false;
    digitalWrite(LED_PIN, LOW);
  }
}
\end{lstlisting}
\end{frame}

\begin{frame}{Por que entender tipos?}
  \begin{itemize}
    \item Arduino tem memória limitada, cada byte importa.
    \item Tipo errado pode causar comportamento inesperado.
    \item Conversão entre tipos é necessária frequentemente.
    \item Variáveis globais devem ser usadas com cuidado em projetos maiores.
  \end{itemize}
\end{frame}

\end{document}
